\section{Conclusion}
\label{sec:conclusion}

Asset context is an under-used, already-collected signal for critical-infrastructure and
public-sector vulnerability remediation. Adding a pre-registered, criticality-dominant context factor to a
hygiene-augmented base scorer~\cite{paper4hygieneprio} delivers a large, mechanism-confirmed
mission-precision gain in the capacity-scarce triage regime, $+9.5$~pp MWP@50 and a near-doubling
of NDCG@50, at no measurable cost to context-blind precision, entirely driven by fleet
heterogeneity, and carried almost entirely by asset criticality. The benefit decays with capacity
and is nothing on a uniform fleet, so we frame CAP-G as a triage-regime tool and report the
capacity-flat primary hypothesis as partially supported rather than re-cutting the bar to
claim a clean win.

The practical message for a public-sector security operations team is narrow and concrete. Where the
remediation window is tight and the fleet is mission-heterogeneous, which is the common operating
condition, ordering the queue by a criticality-weighted context factor moves mission-critical
exposure to the front for the cost of a scoring change over data the organization already maintains;
where capacity is ample or the fleet is uniform, the simpler context-blind scorer is sufficient
and the context layer adds nothing worth its complexity. The ablation sharpens this further: a
criticality-only context layer captures essentially all of the benefit on the target we defined,
so a deployable version of CAP-G can be simpler than its full three-dimensional form. Beyond the
specific scorer, the study is offered as an example of how to bound the value of a prioritization
signal honestly, by pre-registering the bar, isolating the mechanism with an exact control, and
reporting a near-miss as a near-miss. Future work should validate the context mixes against a real
public-sector asset inventory, test whether the criticality-only context layer generalizes across fleet
compositions, and situate the scorer within a continuously self-assessing remediation
loop~\cite{kephart2003vision} where host and exploit state, and therefore the right ordering,
change over time.
